% ============================================================
%  EQUAÇÕES DIFERENCIAIS ORDINÁRIAS DE PRIMEIRA ORDEM
%  Guia de Estudo Completo para Estudantes de Engenharia e Física
% ============================================================
\documentclass[12pt, a4paper]{article}

% ── Packages ────────────────────────────────────────────────
\usepackage[utf8]{inputenc}
\usepackage[T1]{fontenc}
\usepackage[portuguese]{babel}
\usepackage{lmodern}
\usepackage{microtype}

\usepackage[top=2.5cm, bottom=2.5cm, left=2.5cm, right=2.5cm]{geometry}

\usepackage{amsmath, amssymb, amsthm}
\usepackage{mathtools}

\usepackage{graphicx}
\usepackage{tikz}
\usetikzlibrary{arrows.meta, positioning, calc}
\usepackage{pgfplots}
\pgfplotsset{compat=1.18}

\usepackage{booktabs}
\usepackage{array}
\usepackage{tabularx}

\usepackage[dvipsnames]{xcolor}
\definecolor{primary}{RGB}{0,72,144}
\definecolor{accent}{RGB}{220,50,47}
\definecolor{soft}{RGB}{240,246,255}
\definecolor{softgray}{RGB}{245,245,245}
\definecolor{warnyellow}{RGB}{255,248,220}
\definecolor{greenbg}{RGB}{235,255,235}
\definecolor{redbg}{RGB}{255,235,235}

\usepackage[most]{tcolorbox}
\tcbuselibrary{breakable, skins}

\usepackage{enumitem}
\usepackage{multicol}
\usepackage{fancyhdr}
\usepackage{titlesec}
\usepackage{caption}
\usepackage{subcaption}
\usepackage{hyperref}
\usepackage{cleveref}

% ── Configuração do Hyperref ─────────────────────────────────
\hypersetup{
  colorlinks  = true,
  linkcolor   = primary,
  urlcolor    = primary,
  citecolor   = primary,
  pdftitle    = {EDOs de Primeira Ordem -- Guia de Estudo},
  pdfauthor   = {},
  pdfsubject  = {Matemática para Engenharia e Física},
}

% ── Formatação das secções ───────────────────────────────────
\titleformat{\section}
  {\Large\bfseries\color{primary}}
  {\thesection.}{0.5em}{}
  [\color{primary}\titlerule]

\titleformat{\subsection}
  {\large\bfseries\color{primary!80!black}}
  {\thesubsection.}{0.5em}{}

\titleformat{\subsubsection}
  {\normalsize\bfseries\color{primary!60!black}}
  {\thesubsubsection.}{0.5em}{}

% ── Cabeçalho / Rodapé ──────────────────────────────────────
\pagestyle{fancy}
\fancyhf{}
\fancyhead[L]{\small\color{primary}\textbf{EDOs de Primeira Ordem}}
\fancyhead[R]{\small\color{primary}Guia de Estudo}
\fancyfoot[C]{\small\thepage}
\renewcommand{\headrulewidth}{0.4pt}
\renewcommand{\headrule}{%
  \hbox to\headwidth{\color{primary}\leaders\hrule height \headrulewidth\hfill}}

% ── Estilos personalizados de tcolorbox ──────────────────────
\tcbset{
  basestyle/.style={
    breakable, enhanced, fonttitle=\bfseries,
    coltitle=white,
    attach boxed title to top left={yshift=-2mm, xshift=4mm},
    boxed title style={rounded corners, size=small},
  }
}

\newtcolorbox{defbox}[1][]{
  basestyle,
  colback=soft, colframe=primary,
  title={Definição},
  boxed title style={colback=primary},
  #1
}

\newtcolorbox{thmbox}[1][]{
  basestyle,
  colback=greenbg, colframe=green!60!black,
  title={Teorema},
  boxed title style={colback=green!60!black},
  #1
}

\newtcolorbox{exbox}[2][]{
  basestyle,
  colback=softgray, colframe=gray!60,
  title={Exemplo~#2},
  boxed title style={colback=gray!60},
  #1
}

\newtcolorbox{solbox}[1][]{
  breakable, enhanced,
  colback=white, colframe=primary!40,
  title={\textbf{Solução}},
  fonttitle=\bfseries\color{primary},
  borderline west={3pt}{0pt}{primary},
  #1
}

\newtcolorbox{warnbox}[1][]{
  breakable, enhanced,
  colback=warnyellow, colframe=orange!80!black,
  title={\textbf{Erro Comum}},
  fonttitle=\bfseries\color{orange!80!black},
  borderline west={4pt}{0pt}{orange!80!black},
  #1
}

\newtcolorbox{tipbox}[1][]{
  breakable, enhanced,
  colback=greenbg, colframe=green!60!black,
  title={\textbf{Dica}},
  fonttitle=\bfseries\color{green!60!black},
  borderline west={4pt}{0pt}{green!60!black},
  #1
}

\newtcolorbox{exercisebox}[2][]{
  basestyle,
  colback=redbg!40!white, colframe=accent,
  title={Exercício~#2},
  boxed title style={colback=accent},
  #1
}

\newtcolorbox{oralbox}[1][]{
  basestyle,
  colback=primary!8, colframe=primary!70,
  title={Questão de Defesa Oral},
  boxed title style={colback=primary!70},
  #1
}

% ── Ambientes tipo teorema ───────────────────────────────────
\newtheorem{proposition}{Proposição}[section]
\newtheorem{remark}{Observação}[section]

% ── Macros úteis ────────────────────────────────────────────
\newcommand{\R}{\mathbb{R}}
\newcommand{\dd}{\mathrm{d}}
\newcommand{\ODE}{\textsc{edo}}

% ============================================================
\begin{document}

% ── Página de Título ─────────────────────────────────────────
\begin{titlepage}
  \centering
  \vspace*{1.5cm}

  {\color{primary}\rule{\linewidth}{1.5pt}}
  \vspace{0.5cm}

  {\Huge\bfseries\color{primary} Equações Diferenciais Ordinárias\\[4pt]
   de Primeira Ordem}

  \vspace{0.4cm}
  {\color{primary}\rule{\linewidth}{0.5pt}}
  \vspace{0.8cm}

  {\large\itshape Guia de Estudo}

  \vspace{1.2cm}

  \begin{tcolorbox}[colback=soft, colframe=primary, width=0.90\linewidth]
    \centering
    \textbf{Tópicos Abordados}\\[6pt]
    \begin{tabular}{ll}
      \textbullet\ Definição e Classificação          & \textbullet\ Equações Separáveis \\
      \textbullet\ Soluções Gerais e Particulares     & \textbullet\ EDOs Lineares de Primeira Ordem \\
      \textbullet\ Aplicações em Engenharia           & \textbullet\ Estratégias de Verificação \\
      \textbullet\ Exemplos Resolvidos                & \textbullet\ Questões de Defesa Oral \\
    \end{tabular}
  \end{tcolorbox}

  \vspace{1.5cm}

  % Campo de direções para dy/dx = y(1-y), desenhado com duas curvas solução
  \begin{tikzpicture}
    \begin{axis}[
      width=10cm, height=7cm,
      xlabel={$x$}, ylabel={$y$},
      xmin=-0.2, xmax=4.5,
      ymin=-0.3, ymax=1.8,
      axis lines=center,
      tick style={draw=none},
      xticklabels={}, yticklabels={},
      title={\small\color{primary}
             Campo de Direções de $\dfrac{\dd y}{\dd x} = y(1-y)$},
    ]
      % Linhas de equilíbrio
      \addplot[color=primary, thick, domain=0:4.4] {1}
        node[right, font=\tiny] {$y=1$};
      \addplot[color=accent, thick, domain=0:4.4] {0}
        node[right, font=\tiny] {$y=0$};
      % Curva logística por baixo (y0 ~ 0.05)
      \addplot[color=primary!60!black, very thick, domain=0:4.4, samples=80]
        { 1/(1 + 19*exp(-x)) };
      % Solução por cima (y0 = 1.5)
      \addplot[color=gray!70, thick, domain=0:3.0, samples=80]
        { 1/(1 - (1/3)*exp(-x)) };
    \end{axis}
  \end{tikzpicture}

  \vfill
  {\small\color{gray} Matemática para Engenharia e Física\quad|\quad\today}
\end{titlepage}

% ── Índice ───────────────────────────────────────────────────
\tableofcontents
\newpage

% ============================================================
\section{Introdução às Equações Diferenciais}
\label{sec:intro}
% ============================================================

Uma \textbf{equação diferencial} é uma equação que relaciona uma função desconhecida com
uma ou mais das suas derivadas. As equações diferenciais são a linguagem natural das
ciências físicas: descrevem como os sistemas variam no tempo ou no espaço.

\begin{defbox}
  Uma \textbf{equação diferencial ordinária (EDO) de primeira ordem} é uma equação da forma
  \[
    F\!\bigl(x,\, y,\, y'\bigr) = 0,
    \qquad\text{ou equivalentemente}\qquad
    \frac{\dd y}{\dd x} = f(x, y),
  \]
  onde $y = y(x)$ é uma função desconhecida de uma única variável independente $x$, e
  $y' = \dfrac{\dd y}{\dd x}$ é a sua derivada primeira.
\end{defbox}

\subsection{Porque São Importantes as EDOs de Primeira Ordem}

\begin{multicols}{2}
  \textbf{Em física:}
  \begin{itemize}[noitemsep]
    \item Lei do arrefecimento de Newton
    \item Desintegração radioativa
    \item Carga de circuito RC
    \item Escoamento de fluidos (lei de Torricelli)
  \end{itemize}

  \columnbreak

  \textbf{Em engenharia/computação:}
  \begin{itemize}[noitemsep]
    \item Dinâmica populacional
    \item Gradiente descendente em ML
    \item Difusão de calor (1-D)
    \item Cinética de reações químicas
  \end{itemize}
\end{multicols}

% ============================================================
\section{Classificação das EDOs de Primeira Ordem}
\label{sec:classification}
% ============================================================

\subsection{Ordem e Grau}

\begin{defbox}
  \begin{itemize}
    \item A \textbf{ordem} de uma EDO é a ordem da derivada mais elevada que nela aparece.
    \item O \textbf{grau} é o expoente da derivada de maior ordem (após eliminar frações e
          radicais envolvendo a derivada).
  \end{itemize}
\end{defbox}

Todas as equações tratadas neste guia são de \emph{primeira ordem} (a derivada mais elevada
é $y'$) e, salvo indicação em contrário, de \emph{primeiro grau}.

\subsection{Linearidade}

\begin{defbox}
  Uma EDO de primeira ordem é \textbf{linear} se puder ser escrita na forma
  \[
    \frac{\dd y}{\dd x} + P(x)\,y = Q(x),
  \]
  onde $P$ e $Q$ dependem apenas de $x$. É \textbf{não linear} nos restantes casos.
\end{defbox}

\begin{center}
\begin{tabular}{lll}
  \toprule
  \textbf{Equação} & \textbf{Tipo} & \textbf{Motivo} \\
  \midrule
  $y' + 2xy = e^x$  & Linear                      & $P(x)=2x$, $Q(x)=e^x$ \\
  $y' = 3y$          & Linear (coef.\ constantes)  & $P=-3$, $Q=0$ \\
  $y' = y^2$         & Não linear                  & termo $y^2$ \\
  $y' = \sin y$      & Não linear (autónoma)       & $f$ depende apenas de $y$ \\
  $yy' = x$          & Não linear                  & produto $yy'$ \\
  \bottomrule
\end{tabular}
\end{center}

\subsection{Autonomia}

\begin{defbox}
  Uma EDO de primeira ordem é \textbf{autónoma} se $f$ não depender explicitamente de $x$:
  \[
    \frac{\dd y}{\dd x} = f(y).
  \]
\end{defbox}

As equações autónomas admitem \textbf{soluções de equilíbrio (estacionárias)}, obtidas
fazendo $f(y^{*})=0$. O campo de direções é idêntico em todas as retas verticais.

\subsection{Tabela Resumo de Classificação}

\begin{center}
\renewcommand{\arraystretch}{1.3}
\begin{tabularx}{0.92\linewidth}{X X X}
  \toprule
  \textbf{Categoria} & \textbf{Forma Padrão} & \textbf{Característica Principal} \\
  \midrule
  Separável    & $g(y)\,\dd y = h(x)\,\dd x$          & Variáveis desacoplam completamente \\
  Linear       & $y' + P(x)y = Q(x)$                  & Resolvida por fator integrante \\
  Exata        & $M\,\dd x + N\,\dd y = 0$, $M_y=N_x$ & Existe função potencial \\
  Bernoulli    & $y' + P(x)y = Q(x)y^n$               & Linearizada por $v=y^{1-n}$ \\
  Autónoma     & $y' = f(y)$                           & Sem dependência explícita em $x$ \\
  \bottomrule
\end{tabularx}
\end{center}

% ============================================================
\section{Soluções: Geral e Particular}
\label{sec:solutions}
% ============================================================

\begin{defbox}
  \begin{itemize}
    \item A \textbf{solução geral} de uma EDO de primeira ordem contém uma constante
          arbitrária $C$, representando uma família de curvas solução.
    \item A \textbf{solução particular} obtém-se especificando uma
          \textbf{condição inicial} $y(x_0) = y_0$, que determina $C$ de forma única.
    \item Uma \textbf{solução singular} satisfaz a EDO mas não pode ser obtida a partir
          da solução geral por nenhuma escolha de $C$.
  \end{itemize}
\end{defbox}

\begin{thmbox}[title={Existência e Unicidade (Picard--Lindel\"{o}f)}]
  Se $f(x,y)$ e $\dfrac{\partial f}{\partial y}$ forem contínuas num retângulo $R$
  contendo $(x_0, y_0)$, então o problema de valor inicial
  \[
    \frac{\dd y}{\dd x} = f(x,y), \qquad y(x_0)=y_0
  \]
  tem uma solução \textbf{única} em algum intervalo $|x - x_0| < h$.
\end{thmbox}

\begin{warnbox}
  \textbf{Não assuma que a solução existe globalmente.} O intervalo de existência pode ser
  finito. Por exemplo, $y' = y^2$, $y(0)=1$ dá $y = 1/(1-x)$, que diverge em $x=1$.
\end{warnbox}

% ============================================================
\section{Equações Separáveis}
\label{sec:separable}
% ============================================================

\subsection{Definição e Método de Resolução}

\begin{defbox}
  Uma EDO de primeira ordem é \textbf{separável} se puder ser escrita na forma
  \[
    \frac{\dd y}{\dd x} = g(x)\,h(y),
  \]
  isto é, o membro direito fatora num produto de uma função apenas de $x$ e de uma
  função apenas de $y$.
\end{defbox}

\noindent\textbf{Algoritmo para equações separáveis:}

\begin{enumerate}
  \item Reescrever como $\dfrac{\dd y}{h(y)} = g(x)\,\dd x$ \quad (desde que $h(y)\neq 0$).
  \item Integrar ambos os membros:
        $\displaystyle\int \frac{\dd y}{h(y)} = \int g(x)\,\dd x + C$.
  \item Resolver em ordem a $y$ explicitamente se possível; caso contrário, deixar na forma implícita.
  \item Aplicar a condição inicial para determinar $C$ (se fornecida).
  \item \textbf{Verificar soluções de equilíbrio} a partir de $h(y)=0$ --- estas podem
        perder-se ao dividir.
\end{enumerate}

\begin{tipbox}
  Verifique sempre $h(y)=0$ para soluções constantes (de equilíbrio) \emph{antes} de
  separar, pois dividir por $h(y)$ só é válido quando $h(y)\neq 0$.
\end{tipbox}

% ── Exemplo 1 ───────────────────────────────────────────────
\begin{exbox}{1 --- EDO Separável Básica}
  Resolver $\dfrac{\dd y}{\dd x} = xy$.
\end{exbox}

\begin{solbox}
  \textbf{Passo 1 --- Soluções de equilíbrio.} Fazendo $h(y)=y=0$ obtém-se $y\equiv 0$
  (solução válida, facilmente verificada).

  \textbf{Passo 2 --- Separar variáveis} (assumindo $y\neq 0$):
  \[
    \frac{\dd y}{y} = x\,\dd x.
  \]

  \textbf{Passo 3 --- Integrar ambos os membros:}
  \[
    \int \frac{\dd y}{y} = \int x\,\dd x
    \implies \ln|y| = \frac{x^2}{2} + C_1.
  \]

  \textbf{Passo 4 --- Resolver em ordem a $y$:}
  \[
    y = A\,e^{x^2/2}, \quad A\in\R.
  \]
  ($A=0$ recupera a solução de equilíbrio.)

  \textbf{Solução geral:}
  \[
    \boxed{y = A\,e^{x^2/2}, \quad A\in\R.}
  \]
\end{solbox}

% ── Exemplo 2 ───────────────────────────────────────────────
\begin{exbox}{2 --- Problema de Valor Inicial (PVI)}
  Resolver $\dfrac{\dd y}{\dd x} = -2xy^2$, $\;y(0) = 1$.
\end{exbox}

\begin{solbox}
  \textbf{Equilíbrio:} $y=0$ é uma solução.

  \textbf{Separar} ($y\neq 0$):
  \[
    \frac{\dd y}{y^2} = -2x\,\dd x
    \implies -\frac{1}{y} = -x^2 + C
    \implies y = \frac{1}{x^2 - C}.
  \]

  \textbf{Aplicar CI} $y(0)=1$:
  \[
    1 = \frac{1}{0 - C} \implies C = -1.
  \]

  \textbf{Solução particular:}
  \[
    \boxed{y = \frac{1}{x^2+1}.}
  \]

  \textbf{Verificação:}
  $y' = -\dfrac{2x}{(x^2+1)^2}$ e $-2xy^2 = -\dfrac{2x}{(x^2+1)^2}$.
  Ambos os membros coincidem. \checkmark
\end{solbox}

\begin{warnbox}
  \textbf{Erro --- esquecer a constante de integração.} Escrever
  $\ln|y| = x^2/2$ em vez de $\ln|y| = x^2/2 + C$ elimina toda a família de
  soluções, ficando apenas uma curva particular.
\end{warnbox}

% ── Exemplo 3 ───────────────────────────────────────────────
\begin{exbox}{3 --- Solução Implícita}
  Resolver $\dfrac{\dd y}{\dd x} = \dfrac{x^2}{1-y^2}$.
\end{exbox}

\begin{solbox}
  Separar variáveis:
  \[
    (1-y^2)\,\dd y = x^2\,\dd x.
  \]
  Integrar:
  \[
    y - \frac{y^3}{3} = \frac{x^3}{3} + C.
  \]
  Este cúbico em $y$ não tem forma fechada simples; a \textbf{solução geral implícita} é
  \[
    \boxed{y - \frac{y^3}{3} = \frac{x^3}{3} + C.}
  \]
  Uma solução implícita é perfeitamente aceitável --- nem toda a EDO admite forma explícita.
\end{solbox}

% ============================================================
\section{EDOs Lineares de Primeira Ordem}
\label{sec:linear}
% ============================================================

\subsection{Método do Fator Integrante}

A forma padrão é
\[
  y' + P(x)\,y = Q(x).
\]

\begin{thmbox}[title={Fator Integrante}]
  Multiplicar ambos os membros por $\mu(x) = e^{\int P(x)\,\dd x}$. O membro esquerdo
  torna-se uma derivada exata:
  \[
    \frac{\dd}{\dd x}\!\bigl[\mu(x)\,y\bigr] = \mu(x)\,Q(x).
  \]
  Integrando obtém-se
  \[
    y = \frac{1}{\mu(x)}\left[\int \mu(x)\,Q(x)\,\dd x + C\right].
  \]
\end{thmbox}

% ── Exemplo 4 ───────────────────────────────────────────────
\begin{exbox}{4 --- EDO Linear com Fator Integrante}
  Resolver $\dfrac{\dd y}{\dd x} - \dfrac{2}{x}\,y = x^2$.
\end{exbox}

\begin{solbox}
  Aqui $P(x) = -2/x$, $Q(x)=x^2$.

  \textbf{Fator integrante:}
  \[
    \mu = e^{\int -2/x\,\dd x} = e^{-2\ln|x|} = x^{-2}.
  \]

  \textbf{Multiplicar:}
  \[
    x^{-2}y' - 2x^{-3}y = 1
    \implies \frac{\dd}{\dd x}\!\left[\frac{y}{x^2}\right] = 1.
  \]

  \textbf{Integrar:}
  \[
    \frac{y}{x^2} = x + C
    \implies \boxed{y = x^3 + Cx^2.}
  \]

  \textbf{Verificação:} $y'=3x^2+2Cx$. Então
  $y' - \tfrac{2}{x}y = 3x^2+2Cx - 2x^2 - 2Cx = x^2$. \checkmark
\end{solbox}

% ============================================================
\section{Aplicações em Engenharia e Física}
\label{sec:applications}
% ============================================================

\subsection{Lei do Arrefecimento de Newton}

\begin{tcolorbox}[
  colback=soft, colframe=primary,
  title={\textbf{Modelo Físico}},
  fonttitle=\bfseries\color{white},
  attach boxed title to top left={yshift=-2mm, xshift=4mm},
  boxed title style={colback=primary, rounded corners, size=small}
]
  A taxa de variação da temperatura $T$ de um objeto é proporcional à diferença
  entre $T$ e a temperatura ambiente $T_a$:
  \[
    \frac{\dd T}{\dd t} = -k(T - T_a), \quad k > 0.
  \]
\end{tcolorbox}

% ── Exemplo 5 ───────────────────────────────────────────────
\begin{exbox}{5 --- Arrefecimento de Newton (Engenharia)}
  Um componente metálico a $200\,^{\circ}\mathrm{C}$ é colocado num ambiente a
  $25\,^{\circ}\mathrm{C}$. Após 10 minutos arrefeceu até $120\,^{\circ}\mathrm{C}$.
  Determinar $T(t)$ e o instante em que $T = 50\,^{\circ}\mathrm{C}$.
\end{exbox}

\begin{solbox}
  \textbf{Passo 1 --- Formulação.} Seja $u = T - 25$:
  \[
    \frac{\dd u}{\dd t} = -k\,u, \quad u(0)=175.
  \]

  \textbf{Passo 2 --- Resolver.} Separando e integrando:
  $u = 175\,e^{-kt}$, logo $T(t) = 25 + 175\,e^{-kt}$.

  \textbf{Passo 3 --- Determinar $k$.} Usar $T(10)=120$:
  \[
    120 = 25 + 175\,e^{-10k}
    \implies e^{-10k} = \frac{95}{175} = \frac{19}{35}
    \implies k = \frac{1}{10}\ln\!\frac{35}{19} \approx 0{,}061\;\mathrm{min}^{-1}.
  \]

  \textbf{Passo 4 --- Determinar o instante para $T=50\,^{\circ}\mathrm{C}$:}
  \[
    50 = 25 + 175\,e^{-kt}
    \implies e^{-kt} = \frac{1}{7}
    \implies t = \frac{\ln 7}{k}
               = 10\,\frac{\ln 7}{\ln(35/19)} \approx 31{,}8\;\mathrm{min}.
  \]

  \[
    \boxed{T(t) = 25 + 175\,e^{-0{,}061\,t}\;(^{\circ}\mathrm{C}).}
  \]
\end{solbox}

\subsection{Circuito Elétrico RC}

A tensão $V(t)$ aos terminais do condensador num circuito RC em série satisfaz
\[
  RC\,\frac{\dd V}{\dd t} + V = V_s,
\]
onde $V_s$ é a tensão da fonte. Trata-se de uma EDO linear de primeira ordem.

% ── Exemplo 6 ───────────────────────────────────────────────
\begin{exbox}{6 --- Carga de Circuito RC}
  Um circuito tem $R=1\,\mathrm{k}\Omega$, $C=1\,\mathrm{mF}$, $V_s = 5\,\mathrm{V}$,
  e tensão inicial $V(0)=0$. Determinar $V(t)$.
\end{exbox}

\begin{solbox}
  $RC = 1$ s. A EDO torna-se $V' + V = 5$, $V(0)=0$.

  \textbf{Solução homogénea:} $V_h = Ce^{-t}$.

  \textbf{Solução particular:} $V_p = 5$ (constante).

  \textbf{Solução geral:} $V = 5 + Ce^{-t}$.

  \textbf{CI:} $0 = 5 + C \Rightarrow C=-5$.

  \[
    \boxed{V(t) = 5\!\left(1 - e^{-t}\right)\;\mathrm{V}.}
  \]
  A tensão aproxima-se de $V_s = 5\,\mathrm{V}$ exponencialmente com constante de tempo
  $\tau = RC = 1\,\mathrm{s}$.
\end{solbox}

\subsection{Desintegração Radioativa e Datação por Carbono}

\[
  \frac{\dd N}{\dd t} = -\lambda N, \quad N(0)=N_0
  \implies N(t) = N_0\,e^{-\lambda t}.
\]
O \textbf{período de semi-desintegração} é $t_{1/2} = \dfrac{\ln 2}{\lambda}$.

\subsection{Crescimento Populacional Logístico}

O modelo logístico considera uma capacidade de carga $K$:
\[
  \frac{\dd P}{\dd t} = rP\!\left(1-\frac{P}{K}\right).
\]
Trata-se de uma EDO não linear separável (do tipo Bernoulli). A sua solução geral é a sigmoide:
\[
  P(t) = \frac{K}{1 + \!\left(\dfrac{K-P_0}{P_0}\right)\!e^{-rt}}.
\]

\subsection{Difusão de Calor --- Regime Estacionário 1-D}

A temperatura estacionária $T(x)$ ao longo de uma barra com geração interna de calor
satisfaz uma EDO de segunda ordem, mas com condição de fronteira de fluxo prescrito
(condição de Robin) o problema reduz-se a um sistema acoplado de EDOs de primeira ordem
tratável pelos métodos deste capítulo.

\subsection{Gradiente Descendente e Redes Neuronais}

Em aprendizagem automática, a regra de atualização dos parâmetros
\[
  \theta_{k+1} = \theta_k - \eta\,\nabla_{\!\theta}\,\mathcal{L}(\theta_k)
\]
é a discretização de Euler explícito da EDO em tempo contínuo
\[
  \frac{\dd\theta}{\dd t} = -\nabla_{\!\theta}\,\mathcal{L}(\theta),
\]
denominada equação de \emph{fluxo gradiente}. Compreender esta ligação conecta as EDOs de
primeira ordem diretamente aos algoritmos modernos de otimização.

% ============================================================
\section{Estratégias de Verificação e Lista de Erros Comuns}
\label{sec:verification}
% ============================================================

\subsection{Como Verificar uma Solução}

\begin{tipbox}
  \textbf{Protocolo de verificação em três passos:}
  \begin{enumerate}
    \item \textbf{Diferenciar} a solução proposta $y(x)$ para obter $y'(x)$.
    \item \textbf{Substituir} $y$ e $y'$ na EDO original e verificar se ambos os membros
          são iguais para todo $x$.
    \item \textbf{Verificar a condição inicial}: substituir $x=x_0$ e confirmar
          $y(x_0)=y_0$.
  \end{enumerate}
\end{tipbox}

\subsection{Erros Comuns}

\begin{warnbox}
  \textbf{Erro 1 --- Esquecer a constante de integração.}\\
  Escrever $\ln|y| = x^2/2$ em vez de $\ln|y| = x^2/2 + C$ perde toda a família de
  soluções.
\end{warnbox}

\begin{warnbox}
  \textbf{Erro 2 --- Perder soluções de equilíbrio.}\\
  Dividir por $h(y)$ é inválido quando $h(y)=0$. Verificar sempre $h(y)=0$ para soluções
  constantes antes de separar.
\end{warnbox}

\begin{warnbox}
  \textbf{Erro 3 --- Fator integrante incorreto.}\\
  O fator integrante é $\mu = e^{\int P(x)\,\dd x}$, \emph{não} $e^{P(x)}$.
  Esquecer de integrar $P$ é um erro muito frequente.
\end{warnbox}

\begin{warnbox}
  \textbf{Erro 4 --- Aplicar a CI antes de determinar $C$.}\\
  Escrever sempre a solução geral primeiro e só depois aplicar a CI para determinar $C$.
\end{warnbox}

\begin{warnbox}
  \textbf{Erro 5 --- Ignorar o domínio da solução.}\\
  Os logaritmos exigem argumentos positivos; as raízes quadradas exigem argumentos não
  negativos. Indicar sempre o domínio de validade da solução.
\end{warnbox}

\begin{warnbox}
  \textbf{Erro 6 --- Erros de sinal na separação.}\\
  Ao reescrever $f(x,y)=g(x)h(y)$, verificar sempre os sinais, em especial em
  expressões do tipo $-(T-T_a)$.
\end{warnbox}

\subsection{Guia Rápido de Identificação}

\begin{center}
\renewcommand{\arraystretch}{1.4}
\begin{tabular}{p{4.0cm} p{4.0cm} p{5.0cm}}
  \toprule
  \textbf{Reconhecer se\ldots} & \textbf{Tipo} & \textbf{Método} \\
  \midrule
  $y'= g(x)h(y)$                      & Separável  & Separar e integrar \\
  $y' + P(x)y = Q(x)$                 & Linear     & Fator integrante \\
  $M\,\dd x + N\,\dd y=0$, $M_y=N_x$ & Exata      & Encontrar potencial $F$ \\
  $y' + Py = Qy^n$                    & Bernoulli  & Substituir $v=y^{1-n}$ \\
  $y'=f(y)$ apenas                     & Autónoma   & Linha de fase / separar \\
  \bottomrule
\end{tabular}
\end{center}

% ============================================================
\section{Exercícios com Soluções Passo a Passo}
\label{sec:exercises}
% ============================================================

\begin{exercisebox}{1}
  Resolver $\dfrac{\dd y}{\dd x} = \dfrac{y}{x}$, $\;y(1)=3$.
\end{exercisebox}

\begin{solbox}
  Separar: $\dfrac{\dd y}{y} = \dfrac{\dd x}{x}$.

  Integrar: $\ln|y| = \ln|x| + C_1 \Rightarrow y = Ax$, $A\in\R$.

  CI: $3 = A\cdot 1 \Rightarrow A=3$.

  \[
    \boxed{y = 3x.}
  \]
\end{solbox}

\begin{exercisebox}{2}
  Resolver a equação logística
  $\dfrac{\dd P}{\dd t} = 0{,}2\,P\!\left(1 - \dfrac{P}{500}\right)$, $\;P(0) = 50$.
\end{exercisebox}

\begin{solbox}
  Separar ($P\neq 0$, $P\neq 500$):
  \[
    \frac{\dd P}{P(1-P/500)} = 0{,}2\,\dd t.
  \]
  Frações parciais: $\dfrac{1}{P(1-P/500)} = \dfrac{1}{P} + \dfrac{1/500}{1-P/500}$.

  Integrar:
  \[
    \ln|P| - \ln|1-P/500| = 0{,}2t + C.
  \]
  Exponenciar e simplificar:
  \[
    \frac{P}{1-P/500} = Ae^{0{,}2t}
    \implies P = \frac{500}{1+(500/P_0-1)e^{-0{,}2t}}.
  \]
  CI: $P_0=50$, logo $500/50 - 1 = 9$.
  \[
    \boxed{P(t) = \frac{500}{1+9\,e^{-0{,}2t}}.}
  \]
  Quando $t\to\infty$, $P\to 500$ (capacidade de carga).
\end{solbox}

\begin{exercisebox}{3}
  Um objeto de massa $m$ cai sob a ação da gravidade com resistência do ar proporcional à
  velocidade:
  \[
    m\,\frac{\dd v}{\dd t} = mg - kv, \quad v(0)=0.
  \]
  Determinar $v(t)$ e a velocidade terminal.
\end{exercisebox}

\begin{solbox}
  Reescrever como uma EDO linear: $v' + \dfrac{k}{m}\,v = g$.

  Fator integrante: $\mu = e^{kt/m}$.
  \[
    \frac{\dd}{\dd t}\!\left[e^{kt/m}v\right] = g\,e^{kt/m}
    \implies e^{kt/m}v = \frac{mg}{k}e^{kt/m} + C.
  \]

  CI $v(0)=0$: $C = -mg/k$.

  \[
    \boxed{v(t) = \frac{mg}{k}\!\left(1-e^{-kt/m}\right).}
  \]

  Velocidade terminal: $v_{\infty} = \lim_{t\to\infty} v(t) = \dfrac{mg}{k}$.
\end{solbox}

\begin{exercisebox}{4}
  Resolver a equação de Bernoulli $y' + y = y^3$.
\end{exercisebox}

\begin{solbox}
  Aqui $n=3$. Seja $v = y^{1-3} = y^{-2}$. Então $v' = -2y^{-3}y'$.

  Dividir a EDO por $y^3$:
  \[
    y^{-3}y' + y^{-2} = 1
    \implies -\frac{1}{2}v' + v = 1
    \implies v' - 2v = -2.
  \]
  Trata-se de uma equação linear. Fator integrante $\mu = e^{-2x}$:
  \[
    \left(e^{-2x}v\right)' = -2e^{-2x}
    \implies e^{-2x}v = e^{-2x} + C
    \implies v = 1 + Ce^{2x}.
  \]
  Substituir de volta $v=y^{-2}$:
  \[
    \boxed{y^{-2} = 1 + Ce^{2x}, \quad \text{i.e.,}\quad
           y = \pm\frac{1}{\sqrt{1+Ce^{2x}}}.}
  \]
  (Mais a solução trivial $y=0$.)
\end{solbox}

% ============================================================
\section{Questões Sugeridas para Defesa Oral}
\label{sec:oral}
% ============================================================

\begin{oralbox}
  \textbf{Q1.}\ Enunciar o teorema de existência e unicidade para PVI de primeira ordem.
  Dar um exemplo de PVI em que a unicidade \emph{falha} e explicar porquê.

  \medskip
  \textit{Resposta esperada:} Mencionar Picard--Lindel\"of; exemplo de falha pode ser
  $y' = y^{1/3}$, $y(0)=0$ (a derivada parcial $\partial f/\partial y$ é ilimitada
  em $y=0$, pelo que o teorema não se aplica).
\end{oralbox}

\begin{oralbox}
  \textbf{Q2.}\ Explicar o significado físico do fator integrante na EDO linear
  $y' + P(x)y = Q(x)$. Porque é que multiplicar por $\mu = e^{\int P\,\dd x}$ torna
  o membro esquerdo uma derivada exata?

  \medskip
  \textit{Resposta esperada:} Mostrar algebricamente que $(\mu y)' = \mu y' + \mu'y =
  \mu(y'+Py)$ usando $\mu'=\mu P$.
\end{oralbox}

\begin{oralbox}
  \textbf{Q3.}\ Considerar a EDO autónoma $\dfrac{\dd y}{\dd t} = y(2-y)$.
  \begin{enumerate}[label=(\alph*)]
    \item Determinar todas as soluções de equilíbrio.
    \item Determinar a sua estabilidade usando uma análise de linha de fase.
    \item Esboçar várias curvas solução.
  \end{enumerate}
  \textit{Resposta esperada:} Equilíbrios $y=0$ (instável), $y=2$ (estável); as soluções
  que partem entre 0 e 2 crescem até 2; as que estão acima de 2 decrescem até 2.
\end{oralbox}

\begin{oralbox}
  \textbf{Q4.}\ Um tanque contém 100 L de água com 5 kg de sal. Água pura entra a
  2 L/min; a solução bem misturada drena a 2 L/min. Formular e resolver a EDO para a
  quantidade de sal $S(t)$.

  \medskip
  \textit{Resposta esperada:} $S' = -\dfrac{2S}{100} = -0{,}02S$; solução
  $S = 5e^{-0{,}02t}$ kg.
\end{oralbox}

\begin{oralbox}
  \textbf{Q5.}\ Como se relaciona a EDO do gradiente descendente
  $\dfrac{\dd\theta}{\dd t} = -\nabla_{\!\theta}\,\mathcal{L}(\theta)$
  com o método de Euler explícito? O que determina a taxa de convergência?

  \medskip
  \textit{Resposta esperada:} O passo de Euler $\theta_{k+1}=\theta_k -
  \eta\,\nabla\mathcal{L}$ é Euler explícito com tamanho de passo $\eta$; para uma perda
  quadrática a EDO é linear e a solução decai exponencialmente; a taxa de convergência é
  governada pelos valores próprios da Hessiana.
\end{oralbox}

\begin{oralbox}
  \textbf{Q6.}\ Descrever como surgem EDOs de primeira ordem na difusão de calor 1-D e
  explicar que condições de fronteira são necessárias para obter uma solução única.

  \medskip
  \textit{Resposta esperada:} Regime estacionário: $\dfrac{\dd T}{\dd x}=\mathrm{const}$,
  requerendo dois valores de fronteira. Regime transitório: após semi-discretização no
  espaço, cada nó satisfaz $\dfrac{\dd T_i}{\dd t} = \cdots$, uma EDO de primeira ordem
  no tempo que requer uma condição inicial por nó.
\end{oralbox}

% ============================================================
\section{Lista de Verificação do Programa}
\label{sec:checklist}
% ============================================================

\begin{tcolorbox}[
  colback=soft, colframe=primary,
  title={\textbf{Lista de Competências}},
  fonttitle=\bfseries\color{white},
  attach boxed title to top left={yshift=-2mm, xshift=4mm},
  boxed title style={colback=primary, rounded corners, size=small}
]
  Usar esta lista para confirmar a cobertura completa antes da defesa oral.

  \begin{itemize}[label=\textcolor{primary}{$\square$}, leftmargin=1.5em]
    \item Consigo identificar a ordem e a linearidade de qualquer EDO dada.
    \item Consigo determinar se uma EDO de primeira ordem é separável, linear, exata ou
          de Bernoulli.
    \item Consigo resolver equações separáveis e enunciar todas as soluções de equilíbrio.
    \item Consigo aplicar o método do fator integrante a EDOs lineares de primeira ordem.
    \item Consigo resolver um PVI e verificar que a solução satisfaz tanto a equação como
          a condição inicial.
    \item Consigo enunciar o teorema de Picard--Lindel\"of e dar um contraexemplo quando
          as suas hipóteses falham.
    \item Consigo formular e resolver modelos de arrefecimento de Newton, circuito RC e
          desintegração radioativa.
    \item Consigo realizar uma análise de linha de fase para EDOs autónomas e classificar
          os equilíbrios como estáveis, instáveis ou semi-estáveis.
    \item Consigo explicar a ligação entre o gradiente descendente e a EDO de fluxo
          gradiente.
    \item Consigo descrever como surgem EDOs na difusão de calor 1-D e no treino de redes
          neuronais.
  \end{itemize}
\end{tcolorbox}

% ============================================================
\section{Resumo das Fórmulas Principais}
\label{sec:formulas}
% ============================================================

\begin{tcolorbox}[
  colback=softgray, colframe=gray!60,
  title={\textbf{Formulário Essencial}},
  fonttitle=\bfseries,
  attach boxed title to top left={yshift=-2mm, xshift=4mm},
  boxed title style={colback=gray!60, rounded corners, size=small}
]
  \begin{align*}
    &\text{\textbf{Separável:}} \quad
     \int \frac{\dd y}{h(y)} = \int g(x)\,\dd x + C. \\[6pt]
    &\text{\textbf{Linear (FI):}} \quad
     y = \frac{1}{\mu(x)}\!\left[\int \mu(x)\,Q(x)\,\dd x + C\right],
     \quad \mu = e^{\int P(x)\,\dd x}. \\[6pt]
    &\text{\textbf{Bernoulli:}} \quad
     v = y^{1-n} \text{ transforma a equação em forma linear.} \\[6pt]
    &\text{\textbf{Crescimento/decaimento exponencial:}} \quad
     y' = ky \implies y = y_0\,e^{kx}. \\[6pt]
    &\text{\textbf{Logístico:}} \quad
     y' = ry\!\left(1-\tfrac{y}{K}\right) \implies
     y = \dfrac{K}{1+(K/y_0-1)e^{-rt}}. \\[6pt]
    &\text{\textbf{Arrefecimento de Newton:}} \quad
     T(t) = T_a + (T_0-T_a)\,e^{-kt}. \\[6pt]
    &\text{\textbf{Período de semi-desintegração:}} \quad
     t_{1/2} = \dfrac{\ln 2}{\lambda}.
  \end{align*}
\end{tcolorbox}

% ============================================================
\appendix
\section{Exercícios Adicionais (Sem Solução)}
\label{app:extra}
% ============================================================

Os problemas seguintes destinam-se a prática adicional. As respostas são fornecidas abaixo.

\begin{enumerate}
  \item Resolver $\dfrac{\dd y}{\dd x} = (1+y^2)\tan x$, $\;y(0)=1$.
  \item Resolver $y' + \dfrac{y}{x} = \cos x$, $\;y(\pi)=0$.
  \item Um tanque contém 200 L de salmoura com 10 kg de sal. Salmoura de concentração
        0{,}05 kg/L entra a 3 L/min e sai a 3 L/min. Determinar o teor de sal ao fim
        de 40 minutos.
  \item Resolver $y' = (y-1)(y-3)$. Classificar cada equilíbrio como estável ou instável.
  \item Mostrar que $y' = x + y$ não é separável e depois resolvê-la pelo método do
        fator integrante.
\end{enumerate}

\medskip
\noindent\textbf{Respostas:}
\begin{enumerate}
  \item $\arctan y = -\ln|\cos x| + \tfrac{\pi}{4}$, i.e.\
        $y=\tan\!\left(-\ln|\cos x|+\tfrac{\pi}{4}\right)$.
  \item $y = \dfrac{\sin x}{x}$.
  \item Usando $S' = 0{,}15 - \tfrac{3}{200}S$, $S(0)=10$: o equilíbrio é
        $S^{*}=10$ kg; ao fim de 40 min o tanque está próximo do equilíbrio,
        $S(40)\approx 10{,}0$ kg.
  \item Solução geral: $y = \dfrac{3 - Ce^{-2t}}{1 - Ce^{-2t}}$;
        $y=1$ é instável, $y=3$ é estável.
  \item Fator integrante $e^{-x}$; solução geral $y = Ce^x - x - 1$.
\end{enumerate}

\end{document}